\clearpage
\onecolumn

\section{Prompt Templates}
\label{app:prompts}

This section documents the prompts used in our framework. Curly-brace placeholders are filled at run time. Normal-agent and emotion-classifier prompts are shown in neutral boxes; malicious-agent prompts are shown in red boxes to make their adversarial role visually obvious.

% ---- Box styles: only two colors used throughout ----
\definecolor{npfFrame}{HTML}{4A4A4A}
\definecolor{npfBack}{HTML}{FAFAFA}
\definecolor{mpfFrame}{HTML}{8B1A1A}
\definecolor{mpfBack}{HTML}{FDF6F6}

\newtcolorbox{npromptbox}[1]{%
  enhanced,
  colback=npfBack, colframe=npfFrame,
  boxrule=0.5pt, arc=2pt,
  fontupper=\small,
  fonttitle=\bfseries\small,
  coltitle=white, colbacktitle=npfFrame,
  left=6pt, right=6pt, top=5pt, bottom=5pt,
  before skip=8pt, after skip=8pt,
  title={#1}
}

\newtcolorbox{mpromptbox}[1]{%
  enhanced,
  colback=mpfBack, colframe=mpfFrame,
  boxrule=0.5pt, arc=2pt,
  fontupper=\small,
  fonttitle=\bfseries\small,
  coltitle=white, colbacktitle=mpfFrame,
  left=6pt, right=6pt, top=5pt, bottom=5pt,
  before skip=8pt, after skip=8pt,
  title={#1}
}

\subsection{Normal Agent Prompts}

The system prompt instantiates the agent's identity, psychological profile, the current topic, and the stance scale. Numeric belief dimensions $D_1$--$D_5$ are mapped to natural-language descriptions through the rule in Table~\ref{tab:personality_mapping}.

\begin{npromptbox}{Normal agent: system prompt}
\begin{verbatim}
You are {nickname}, participating in an online forum discussion.

Your personality and beliefs:
{personality}
{intervention_belief}

Topic: {topic_description}

Stance scale (-5 to +5):
  -5    : {stance_negative} (absolute, unwavering)
  -3/-4 : Lean against, with concerns
  -1/-2 : Slight reservations, mostly uncertain
   0    : Genuinely undecided / no opinion
  +1/+2 : Slight support, with caveats
  +3/+4 : Lean in favor, fairly confident
  +5    : {stance_positive} (absolute, unwavering)

Most people hold moderate views (+/-1 to +/-3); extreme
positions (+/-4, +/-5) require strong conviction.
Your current position:   {current_stance}
Your current confidence: {current_confidence}/10

How you participate:
  - Natural and direct, like a real person in an online
    discussion. No literary metaphors.
  - Conversational; incomplete sentences, filler words,
    and casual grammar are all fine.
  - Sometimes respond to others, sometimes share an
    independent thought, sometimes stay quiet when you
    have nothing new to add. Do what feels natural.

You can change your mind if genuinely persuaded by new
information or arguments, but do not flip without reason.
\end{verbatim}
\end{npromptbox}

\begin{table}[!htbp]
\centering
\small
\renewcommand{\arraystretch}{1.15}
\caption{Rule-based mapping from numeric belief dimensions to natural-language personality phrases. Each agent's personality block concatenates one phrase per dimension.}
\label{tab:personality_mapping}
\setlength{\tabcolsep}{4pt}
\begin{tabular}{@{}l p{0.26\textwidth} p{0.26\textwidth} p{0.26\textwidth}@{}}
\toprule
\textbf{Dimension} & \textbf{High ($\geq 7$)} & \textbf{Medium ($4$--$6$)} & \textbf{Low ($\leq 3$)} \\
\midrule
$D_1$ Trust & generally trust institutions, officials, and mainstream media & mixed feelings about institutions, sometimes trusting, sometimes sceptical & distrust institutions, officials, and mainstream narratives \\
$D_2$ Care & care deeply about fairness and the suffering of others & care about fairness but try to stay practical & practical-minded, not swayed by emotional appeals \\
$D_3$ Loyalty & strongly value tradition, loyalty, and social order & respect tradition but open to change & question traditions, prefer individual freedom over group loyalty \\
$D_4$ Just-World & believe the world is generally fair and people get what they deserve & life is sometimes fair, sometimes not & system is often unfair and stacked against ordinary people \\
$D_5$ Analytical & rely heavily on logic and evidence when forming opinions & use a mix of logic and gut feeling & trust gut and personal experience more than abstract data \\
\bottomrule
\end{tabular}
\end{table}

The user prompt drives the BDI reasoning protocol. The agent first writes an internal chain-of-thought between \texttt{[COT]} markers, then emits a JSON object that encodes its action, optional reply target, updated stance and confidence, and the message itself.

\begin{npromptbox}{Normal agent: user prompt}
\begin{verbatim}
Round {round_num}
{own_history}

What others said:
{neighbor_messages}
{intervention_section}

Think briefly, then respond.

[COT]
Beliefs:    <what you noticed from others>
Desires:    <what you want to express right now>
Intentions: <speak / reply / silent, and why>
[/COT]

{
  "action":     "speak" | "reply" | "silent",
  "target":     null    | "<Nickname>",
  "stance":     <int -5..5>,
  "confidence": <int 1..10>,
  "message":    "<2-4 sentences; empty if silent>"
}
\end{verbatim}
\end{npromptbox}

\subsection{Malicious Agent Prompts}

Each malicious agent receives a strategy-specific system prompt that defines what it should do and what it must \emph{not} do, preventing cross-strategy contamination. The user prompt provides the rolling discussion context, the dominant sentiment, and (when applicable) what teammate attackers have already posted in the current round.

\begin{mpromptbox}{Malicious agent: system prompt}
\begin{verbatim}
You are an influence operator posing as "{nickname}" in a
forum discussion.

Your mission: shift the group toward a negative stance on
the topic and destroy trust in the mainstream position.
You must NEVER break character; you are a real person.

Topic: {topic_description}

Your strategy -- follow this STRICTLY:
{strategy_description}

{teammate_instruction}

Core rules:
  1. Stay in character as your strategy type. If pseudo-
     rational, never use emotional stories. If conspiracy,
     talk about systems, not individuals.
  2. Directly engage with what others said -- quote them
     by name, challenge them.
  3. Escalate over time. Early: concerned citizen. Later:
     angry, accusatory.
  4. Write 4-6 sentences. Be substantive. Sound like a
     real person, not an AI.
\end{verbatim}
\end{mpromptbox}

\begin{mpromptbox}{Malicious agent: user prompt}
\begin{verbatim}
Round {round_num} / {total_rounds}

What people are saying:
{sampled_messages}

Context:
  Group mood:       {dominant_sentiment}
  Key debates:      {key_points}
  Your prior posts: {own_previous}

[MAL-COT]
Situation: <group mood, who is vulnerable>
Play:      <your angle this round, strictly your strategy>
[/MAL-COT]

{
  "message": "<post as {nickname}, 4-6 sentences,
               strictly your strategy, engage by name>"
}
\end{verbatim}
\end{mpromptbox}

\subsection{Emotion Classification Prompt}

Per-message emotion labels in Section~\ref{sec:experiments} are produced by a separate LLM-as-classifier pass that uses a few-shot calibration prompt. The full system prompt enforces a strong neutral-by-default rule with eight calibrated examples spanning the Plutchik categories~\citep{Plutchik1980,Mohammad2018}; we summarize the rule set and show the user prompt verbatim.

\begin{npromptbox}{Emotion classifier (system + user)}
\textbf{System.} Classify the emotional tone of an online discussion message into one of \emph{anger, fear, sadness, disgust, trust, anticipation, surprise, neutral}. Each category is defined with explicit calibration rules. Stating that something is wrong or bad is \emph{neutral}, not anger or disgust. Mild concern is \emph{neutral}, not fear. Firm disagreement without insults is \emph{neutral}, not anger. Visceral revulsion language is required for \emph{disgust}. Eight short examples illustrate each category.

\medskip
\textbf{User.}
\begin{verbatim}
Classify:
  Message: "{message}"

Respond with ONLY JSON:
{
  "category": "anger" | "fear" | "sadness" | "disgust"
              | "trust" | "anticipation" | "surprise"
              | "neutral",
  "reason":   "<brief>"
}
\end{verbatim}
\end{npromptbox}
